\section{Wave propagation in homogeneous layers}

In the previous section, a formulation for rigid-frame porous media and two formulations for poroelastic media modelling have been presented. Concluding the presentation of these motion equations, some characteristics of the wave propagation in unbounded porous and poroelastic media have been shown. The present section proposes the introduction of a geometry, that is to say boundaries or interface conditions delimiting different domains characterized each by a different type of material. Generally, the configurations studied in this work assume a monochromatic plane wave originated from a fluid half-space, i.e. air, impinging at the surface of a structure and is then scattered (reflected and transmitted) from and through each interface. The existence of an interface leads to the propagation of not only the bulk waves but also of so-called surface waves \cite{auld1973} where the interface itself acts as the support for propagation. Different types of waves propagate along the interface. An example can be surface water waves that follow the interface between air and water; all the energy is localized at the interface. A similar behaviour exist for seismic waves where the interface of a solid substrate supports the wave propagation, that is Rayleigh (longitudinal motion) and Love waves (transverse motion) which both show an energy decaying with depth. Scholte waves are waves that propagate along the interface between a fluid and a solid medium.

If the material in which the wave propagates is bounded by two parallel, free, and flat surfaces, i.e., a plate, then it is refered to as a layer. The waves supported by the medium propagate by reflecting alternatively from one surface to the other; they are said to be guided along the layer. These observations of successive reflections and refractions lead to the concept of waveguides, but their analytical descriptions often leads to writing complex system that can be difficult to solve. To determine the modes of a structure, one must seek to solve the propagation equations and the boundary conditions imposed at its boundaries. In a first section, the scattering of a plane wave by a homogeneous porous layer is investigated along with the various method available. Next, the methods for the computation of the modes of these configurations are presented. In a subsequent section, more complex geometries, those of metaporous structures, are of interest and both methods for the computation of the scattering coefficients and the modes of those configurations will be discussed.

\subsection{Analytical plane wave modelling}
\label{ssec:GM}

Modeling acoustic propagation in homogeneous layers by assuming plane waves is the simplest formalism in acoustics. We assume these plane waves such that they satisfy the motion equations, e.g. the ones introduced in \cref{eq:chap1_helmholtz} or \cref{eq:chap1_biot}. Introducing a simple problem, that of the scattering of an incident plate by a porous layer, we present an example of plane wave modelling for a simple system. The system is 2D with coordinates $\mbf x = (x_1, x_2)$ constituted of a porous layer with thickness $h$ and thus two interfaces of infinite extent along $x_1$, there is a fluid-porous interface at $x_1 = h$ and a porous-fluid interface at $x_1 = 0$. Both media of the air half-spaces and porous layer are modeled as fluid, thus the interfaces conditions should impose continuity of the pressure fields and normal particle displacements. 

Assuming a $ e^{\im \omega t}$ time-dependance, the incident pressure field originated from the upper domain that satisfy the Helmholtz equation takes the form of, 
\begin{equation}
    p_{0+} = a_+ e^{-ik_1 x_1} e^{-ik_{2,0} (x_2 -h)} + a_- e^{-ik_1 x_1} e^{ik_{2,0} (x_2 -h)}.
\end{equation}
The pressure field in the porous layer is expressed as a sum of forward and backward waves,
\begin{equation}
    p_{p} = b_+ e^{-ik_1 x_1} e^{-ik_{2,eq} x_2} + b_- e^{-ik_1 x_1} e^{ik_2 (x_2-h)},
\end{equation}
and that in the lower air domain is,
\begin{equation}
    p_{0-} = c e^{-ik_1x_1}e^{ik_{2,0} x_2},
\end{equation}
with wavenumbers such that they hold for the Snell-Descartes law that states $k_1^2 + k_{2,0}^2 = k_0^2$ and $k_1^2 + k_{2,eq}^2 = k_{eq}^2$. The wavenumber should be imposed with $k_1 = k_0 \sin{\theta}, k_2 = -k_0 \cos{\theta}$, with $\theta$ the angle of incidence, that governs the direction of propagation of the incident field. The displacement fields are additionally expressed using the momentum equation relating the latter with the pressure fields as $\partial_2 p = -\im \omega \rho_i u_2,\:\: i=$. 

Next, the interface conditions describe the relationships between properties at the boundary between two media, including the continuity or vanishing of given fields. If a media is terminated by a surface (e.g. a rigid wall), a constraint should be applied on the fields to represent it, denoted as boundary conditions. Assuming a continuity between the two fluids in the present case with interfaces located at $x_2 = 0$ and $x_2 = h$, we obtain the following conditions,
\begin{equation}
    p^t(h) = p^l(h), \quad u_{2}^t(h) = u_2^l(h) \quad p^l(0) = p^b(0) \quad _{2}^l(0) = u_2^b(0),
\end{equation}
with superscripts $b$ and $t$ denoting respectively the bottom and top interfaces. 

The full analytical description of a system constituted from one or several interfaces can be done by writing a linear system built from the interface conditions, which vector of unknowns are the amplitudes of each of the existing waves in the system, that is $\begin{pmatrix} b_+ & b_- & a_+ & c \end{pmatrix}^\transp$. The forcing right-hand side corresponds to the incident wave terms. The solution of the system yields all unknown amplitudes of the system, in particular the reflection and transmission coefficients. Imposing the interface conditions leads to writing the following matrix system, 
\begin{equation}
    \begin{pmatrix}
        1 & e^{\im k_{2,eq}h}& 0 & -1 \\
        \cfrac{k_{2,eq}}{\rho_{eq}}& -\cfrac{k_{2,eq}}{\rho_{eq}} e^{\im k_{2,eq}h}&   0& \frac{k_{2,0}}{\rho_0} \\
        e^{\im k_{2,eq}h}& 1 & -1 & 0 \\
        \frac{k_{2,eq}}{\rho_{eq}} e^{\im k_{2,eq}h} &    -\frac{k_{2,eq}}{\rho_{eq}}& -\frac{k_{2,0}}{\rho_{0}}&  0 
    \end{pmatrix} \begin{pmatrix} b_+ \\ b_- \\ a_+ \\ c \end{pmatrix} = \begin{pmatrix}
        0 \\ 0 \\ a_- \\ \frac{k_{2,eq}}{\rho_{eq}} a_-
    \end{pmatrix}. 
    \label{eq:gm_pres_layer}
\end{equation}
\comment{vérifier l'écriture du système.}

From this expression, the system can simply be solved to obtain all scattered amplitudes $b_+, b_-, a_+$ and $c$. In the case of a rigidly-backed layer, the interface at $x_2=h$ is replaced by a rigid termination. For (equivalent) fluids, a hard wall or Neumann-type boundary condition is translated by a vanishing normal particle displacement field along the boundary, that is $u_2(\Gamma) = 0$. The term depending on $c$ thus vanishes, since $p_{0-}$ does not exist; this reduces the size of the system. In this case, the reflection coefficient $R$ corresponds to the amplitude $a_+$.

A few results for the simple cases of the scattering of a plane wave by a rigidly-backed layer and in a symmetric configurations (with air coupling at both interfaces) are represented by the diagrams in Figs~\ref{fig:poroelastic_scattering_ex}a and b and results in Figs~\ref{fig:poroelastic_scattering_ex}c and d, respectively. The GM results are presented by the continuous lines. There, the layer for both porous and PEM are both $h=5$ cm thick and the incident wave propagates with a $\theta=30^o$ angle. The parameters for each material corresponds to the one used in Table~\ref{tab:poroelastic_paramsJAP} from Chapter~\ref{chap:JAP}. The system used to obtain the reflection coefficient for the PEM layer is done in a similar fashion as the equivalent fluid and can be read in the appendices of \cite{weisser2016}.

The first case, also referenced to as a reflection problem compares the equivalent fluid and the full Biot model response. The two curves follow the same trend but additional effects due to the skeleton motion can be observed on the continuous lines while the equivalent fluid response shows to be smoother. For all configurations, the melamine layer is acoustically transparent at lower frequency and absorption can be noticed at slightly higher frequencies. 

Generally, the application of the Global Method allows to write the propagation of plane waves through any stack of layer, by putting together a matrix system that contains all interface conditions in the system, in the general form,
\begin{equation}
    \mbf M(\omega) \mbf x = \mbf f,
    \label{eq:sys_gm}
\end{equation}


    \subsection{Transfer Matrix Method}
The widely used method we know today as the \acrfull{tmm}, was initially developed by the works of Thomson \cite{thomson1950} and Haskell \cite{haskell1953}. 

in order to describe elastic wave propagation in multilayer systems. This method can be applied to a range of wave physics among geophysics, quantum mechanics, electromagnetics and obviously acoustics. This method allows for an easy multiphysics coupling with interface matrices leads to large systems. 

The core of the method is in the evaluation of the state vector that regroups the variables of interest in the system and evaluate them at each interface. propagate the properties stored in the state vector from one layer to the next. An amount of $N$ layers of the same nature are easy to manipulate since $\mbf T_{\mathrm{global}} = \prod_{n=0}^N \mbf T_n$. This method was deep dived in several textbooks among which we propose Refs.~\cite{allard2009, jimenez2021}. It allows for the computation of complex system assembly of elements by introducing additional elements. \cite{doutres2015}. While very versatile, it remains unstable in limit cases at high frequencies or for large thickness of the layers. While the formulation is mathematically exact, the evaluation of the system involves involves finite-precision calculation which leads the very sensitive exponential terms to diverge in the previously specified regimes.
\begin{equation}
    \begin{pmatrix}
        p(0) \\ u_2(0)
    \end{pmatrix} = 
    \begin{pmatrix}
        \cos(k_{eq}h) & \im Z_{eq} \sin(k_{eq}h) \\
        \im Z_{eq}^{-1} \sin(k_{eq}h) & \cos(k_{eq}h) \\
    \end{pmatrix}
    \begin{pmatrix}
        p(h) \\ u_2(h)
    \end{pmatrix}
    \end{equation}
From the evaluation of the transfer matrix of the global system, the scattering matrix can be obtained and thus the reflection and transmission coefficients. 

The absorption coefficient for the same two configurations as the last part is presented in Figs~\ref{fig:poroelastic_scattering_ex}c and d with the markers. It can be seen that the same results are obtained. 

\begin{figure}[h!]
    \centering 
    \begin{tikzpicture}
\def\height{1} \def\pad{.3} \def\arrow{.7}\def\totalen{7}
\begin{scope}[xshift=-10cm, yshift=1.8cm,scale=.5]
    \fill [
        pattern={Lines[distance=2mm,angle=45,line width=0.7mm]},pattern color=black
       ] (0, \height) rectangle ++(\totalen, .5*\pad);
    \draw[top color = fluid, bottom color = white,draw=none] (0, 0) rectangle ++ (\totalen, -6*\pad);
    \draw[fill=pem, draw=none] (0, 0) rectangle ++ (\totalen, \height);

    \draw[black] (0, 0) -- ++ (\totalen, 0);
    \draw[black] (0, \height) -- ++ (\totalen, 0);
    \draw[<->] (-\pad, 0) -- ++ (0, \height) node[midway, left] {$h$};
    \draw[thick, <-] (\totalen/2, 0) -- ++ (210:4*\pad) node[left] {$p_i$};
    \draw[thick, ->] (\totalen/2, 0) -- ++ (-30:4*\pad) node[right] {$R$};
    \draw[dashed] (\totalen/2, 0) -- ++ (0, -4*\pad);
    \draw[<->] (\totalen/2, -2*\pad) arc (90:30:-2*\pad) node[midway, below] {$\theta$};
    
    \node at (-.7, 1.8){a)};
\end{scope}
\begin{scope}[xshift=-10cm, yshift=-1.5cm, scale=.5]
    \draw[top color = fluid, bottom color = white,draw=none] (0, 0) rectangle ++ (\totalen, -6*\pad);
    \draw[top color = white, bottom color = fluid,draw=none] (0, \height) rectangle ++ (\totalen, 6*\pad);
    \draw[fill=pem, draw=none] (0, 0) rectangle ++ (\totalen, \height);

    \draw[black] (0, 0) -- ++ (\totalen, 0);
    \draw[black] (0, \height) -- ++ (\totalen, 0);
    \draw[<->] (-\pad, 0) -- ++ (0, \height) node[midway, left] {$h$};
    \draw[thick, <-] (\totalen/2, 0) -- ++ (210:4*\pad) node[left] {$p_i$};
    \draw[thick, ->] (\totalen/2, 0) -- ++ (-30:4*\pad) node[right] {$R$};
    \draw[thick, ->] (\totalen/2, \height) -- ++ (30:4*\pad) node[right] {$T$};
    \draw[dashed] (\totalen/2, 0) -- ++ (0, -4*\pad);
    \draw[<->] (\totalen/2, -2*\pad) arc (90:30:-2*\pad) node[midway, below] {$\theta$};
    \node at (-.7, 3){b)};
\end{scope}

\begin{scope} 
    \node[draw=none,fill=none] at (0, 0){\includegraphics{poro_elastic_scattering.pdf}};
    \node at (-5.8, 1){c)};   
    \node at (-0.1, 1){d)};   
\end{scope}
\end{tikzpicture}
    \caption[Diagram with results of the reflection and transmission coefficients for equivalent fluid and poroelastic layers]{Diagram for the a) reflection, b) transmission problem. Scattering coefficients from a poroelastic layer (continuous) and equivalent fluid layer (dashed) computed using the GM for the c) reflection, d) transmission problem. Results computed with the TMM are also shown for the absorption coefficient for the equivalent fluid (\smallcirclemarker)  and for the absorption coefficient (\smalltrianglemarker)}
    \label{fig:poroelastic_scattering_ex}
\end{figure}

    \subsection{Guided waves for dispersion problems}
From the previously obtained linear systems in Section~\ref{ssec:GM} for the configurations, and removing the incident wave from the system, i.e. the forcing (right-hand side) term, an homogeneous linear system is obtained, representing the guiding of waves in the layer. This set of equations becomes singular when its determinant vanishes for a given parameter, allowing to identify the modes of the structure. Several formalism can be adopted: solving for frequencies while imposing the wavenumber or the other way around. The first one is useful when interested in the time-dependant effects the structure has on the propagation. In this work, we are interested in the losses generated by the structure and the material and as such frequency is fixed in order to solve wavenumbers. 

        \subsubsection{Computing the modes with semi-analytical tools}

for porous/poroelastic materials, root-finding \cite{boeckx2005} + groby2008
single homogeneous layers

Once the linear system is obtained, we are faced with either or two of these possibilities. Solving wavenumbers can be done through the writing and solving of a (possibly generalized) eigenvalue problem or by root-finding of solutions in the case where writing the first one appears impossible. 

Using Eq.~(\ref{eq:sys_gm}), root finding such that $\mathrm{det}\left(\mbf M(k_1, \omega)\right) = 0$; 
the determinant of the matrix from Eq.~\ref{eq:gm_pres_layer} can be expressed as 
\begin{equation}
    2 \frac{k_{2,eq} k_{2,0}}{\rho_{eq} \rho_{0}} \text{cos}(k_{2,eq} h) + \im \left(\frac{k_{2,eq}^2}{\rho_{eq}^2} + \frac{k_{2,0}^2}{\rho_{0}^2}\right) \text{sin}(k_{2,eq} h) = 0
\end{equation}
        \subsubsection{Numerics for mode computations}

Regarding the computation of complex modes in homogeneous layer, several numerical methods have been investigated in the last decades. Guided waves represents a large topic and we focus here on the case of simple geometries; similar configurations are considered, an unbounded single (or stack of) layers with two infinite interfaces as presented in the section where we were interested in their acoustic response. 

Spectral methods have been around for quite some time and have been used for many numerical applications. We focus for this work on \acrfull{scm}, which has been especially used for the computation of modes for guided waves in layers. \cite{trefethen2000}

In the case of simple 2D geometries, a 1D discretization along the transverse direction associated to a harmonic motion in the longitudinal direction. Even if the use of finite elements seems farfetched, this method could be relevant in contexts where boundary effects or guide with varying cross-sections are studied. However, a variation of finite elements based on a 1D discretization of the waveguide is entitled \acrfull{safe} and is akin to \acrshort{scm} which we have described in the previous paragraph. The 1D discretization is done along the thickness of the layer by regular nodes, and the fields are interpolated using finite elements shape functions. Examples of \acrfull{safe} use cases are very similar to \acrshort{scm} and include \comment{lister des publis ici}. Examples are also given in Chapter \ref{chap:jap}. A detailed description of the calculation of Lamb modes using a simple SAFE implementation using high-order Lobatto polynomials is proposed in Appendix.~\ref{app:safe_vs_scm}. There, its performance is also compared to that of SCM. 

 discuss surface and guided waves in a paragraph

    \subsubsection{A bestiary of the wave solutions in dissipative and open waveguides}
In conservative and closed systems, propagating waves are either guided or trapped waves, while in open non-dissipative systems, the acoustic energy is able to radiate outside of the waveguide. Part of propagating waves are thus attenuated and radiate energy to the open domain. They are thus considered quasi-guided since they are propagating waves with some attenuation. Trapped modes can still occur. New family of modes appear as leaky backward and leaky forward waves that exhibit unique properties. Their counterparts are incoming waves.

The addition of dissipation leads to the mixing of several loss mechanisms: the leakage of waves and the dissipations within the guide (which can be of whatever nature). In this case, the separation between incoming and leaky waves is not possible anymore. This particular aspect has been very scarcely tackled in the literature: an unpublished paper is interested in an elastic waveguide \cite{kiefer2020} while in optics, a paper discusses the waves guided in nanospheres with a complex permittivity \cite{campione2011}\comment{vérifier la citation}.

There is a discussion to be made on the relevance of the imposition of a local coupling contrary to a full modelling of the surrounding media with some termination condition satisfying the Sommerfeld condition, PML or analogous. 

while leading to no assumptions on the radiating field nature contrary to imposing a local coupling, the modeling of the half-space by using some artificial absorbing layer at the end leads to the apparition of modes of the cavities \cite{gallezot2017}
    \subsection{Numerics for the computation of dispersion relations}

\begin{itemize}
    \item presentation of scm, safe + possible with fem
    \item cite chap. 2
\end{itemize}

\begin{figure}
    \centering
    \caption{Discretization for each method}
\end{figure}